\section{Divide and Conquer}

\begin{definition}

\textbf{Divide and Conquer } refers to the strategy of splitting a problem into smaller subproblems,
                solving them recursively, and combining the results.
            
\end{definition}

            This concept can be used in many parallel algorithms. Until now we have used manual fork (.start()) and join
            (.join()) of Threads. We can continue using manual fork/join with D \&C, but this can lead to several
            Problems.

            \begin{itemize}
    \item Poor reusability of code
    \item We need to commit to a number of threads a priori and can not dynamically adjust.
    \item Subproblems are not automatically assigned to a thread.
\end{itemize}
.

            \begin{remark}

                There are some fixes that can improve manual fork/join such as sequential-cutoffs or running one
                subproblem on the thread itself, instad of creating two new threads.
            
\end{remark}

\subsection{Executor Service}

            The Executor Service framework provides the way to schedule and run tasks on a thread-pool automatically.

            \begin{remark}

                Executor service can handle Runnables as well as Callables, with the diffrence that Callables return a
                result.
            
\end{remark}

            Pros and Cons of Executor service:
            \begin{itemize}
    \item Not suited for: Recursive subproblems that depend on each other and have a deep fork-tree. As the
                    service will eventually run out of threads.
    \item Suited for: flat structures that can run independently in parallel.
\end{itemize}

\subsection{Fork-Join}

            The fork-join framework is built for divide and conquer parallelism. It automatically manages a pool of
            threads to solve the recursive tasks using a work-stealing strategy, where threads have local task-queues
            which they work an and if empty steal from other threads.

            \begin{codeexample}

                class myTask extends RecursiveTask {
                myTask(int start, int end) {
                //doWork
                }
                protected compute() {
                if (end-start == 1) { return baseCase();} } else { left=new myTask(start, (start+end)/2);
                right=new myTask((start+end)/2, end); left.fork(); right.fork(); return left.join() + right.join();
                } }
            
\end{codeexample}

\begin{important}

                Make sure that every task has enough work to do (aprox. 100-10'000 ops.), else the overhead will be
                larger than the parallel
                speedup gain.
            
\end{important}

\subsection{Task graphs}

            The execution of parallel programs using Fork-Join can be represented as a Task-Graph (DAG).

            \begin{definition}

                We introduce a new size $T_{\infty}$ which is the execution time assuming unlimited hardware. This is
                called the \textbf{span} and can be visualized as the sum of the execution times of the longest dependent
                path in our
                DAG.
            
\end{definition}

            It is the programmers responsibility to find the optimal decomposition/alogrithm for the problem to ensure a
            minimal span. Fork-join will guarantee optimal scheduling of tha tasks with minimal idletime.

        